\sgaStableTarget{sga6.e05.para.3-3}3.3.
Let @@M0@@, @@M1@@ be two @@M2@@-rings, and @@M3@@ a ring homomorphism. By Definition 3.2.ii), @@M4@@ is a @@M5@@-homomorphism if and only if, for every @@M6@@, and every @@M7@@, @@M8@@, or equivalently @@M9@@.

\sgaStableTarget{sga6.e05.para.3-4}3.4. From relation 3.2.ii), we deduce the relations:
\sgaStableTarget{sga6.e05.eq.3-4-1}@@M10@@
They show that the homomorphism @@M11@@ defines on @@M12@@ a new pre-@@M13@@-ring structure; we shall see that in general it is not a @@M14@@-ring structure. Denoting by @@M15@@ the substitution in the ring of formal series in @@M16@@ with coefficients in @@M17@@ defined by @@M18@@, we have:
@@M19@@
The automorphism @@M20@@ is functorial in @@M21@@, and defines an automorphism of the functor @@M22@@ considered as a functor to abelian groups; with this automorphism, by transport of structure, we can define on @@M23@@ a new structure as a functor to @@M24@@-rings, such that for every @@M25@@-ring @@M26@@, the homomorphism @@M27@@ is a @@M28@@-homomorphism; we shall make explicit the multiplication and the operations @@M29@@ of this second structure.

\sgaStableTarget{sga6.e05.para.3-5}3.5. The multiplication is defined by the homomorphism of functors:
@@M30@@
where the sign @@M31@@ denotes the multiplication defined in 2.3. This new multiplication will be denoted by @@M32@@. As usual, it is defined by its value on @@M33@@; we therefore have:
